\clearpage
\section{МОДЕЛЬ БАНКІВСЬКОГО ПЕРЕКАЗУ ТА ОБРОБЛЕННЯ ВІДМОВ}

\paragraphheading{Послідовність виконання транзакції.}
У моделі виконується переказ суми \(250.00~\mathrm{UAH}\) із рахунку \texttt{101}, розміщеного на \texttt{Node A}, до рахунку \texttt{202}, розміщеного на \texttt{Node B}. Процес оброблення операції складається з таких кроків:
\begin{enumerate}
    \item Координатор створює глобальну операцію переказу \(250.00~\mathrm{UAH}\).
    \item \texttt{Node A} відкриває транзакцію, перевіряє активність рахунку та достатність залишку, виконує попереднє списання й переходить до стану \texttt{PREPARED}.
    \item \texttt{Node B} відкриває транзакцію, перевіряє активність рахунку отримувача, виконує попереднє зарахування й переходить до стану \texttt{PREPARED}.
    \item Якщо обидва вузли відповіли \texttt{READY}, координатор записує рішення \texttt{COMMIT} та виконує \texttt{COMMIT PREPARED} на обох вузлах.
    \item Якщо хоча б один учасник не готовий до фіксації, координатор приймає рішення \texttt{ROLLBACK} і скасовує вже підготовлені зміни.
\end{enumerate}

\paragraphheading{Обґрунтування вибору протоколу.}
Для задачі використано двофазний протокол фіксації (2PC), оскільки банківський переказ не допускає часткового результату: списання без зарахування або зарахування без списання. PostgreSQL надає команди \texttt{PREPARE TRANSACTION}, \texttt{COMMIT PREPARED} та \texttt{ROLLBACK PREPARED}, які дають змогу змоделювати дві фази завершення транзакції.

Роль координатора полягає у створенні ідентифікатора глобальної операції, надсиланні запиту на підготовку вузлам, збиранні відповідей, збереженні остаточного рішення й передаванні його кожному учаснику. Якщо після запису рішення відбувається збій, координатор після відновлення повторює завершальну команду відповідно до журналу рішень.

\paragraphheading{Вплив мережевих затримок і відмов.}
Мережеві затримки збільшують час очікування повідомлень \texttt{READY}/\texttt{ABORT} та період утримання блокувань. Якщо рішення \texttt{COMMIT} ще не було зафіксовано, відсутність відповіді від учасника може призвести до рішення про відкат. Якщо ж рішення \texttt{COMMIT} уже записане, його необхідно повторно передати вузлам після відновлення з'єднання. Можливі ситуації відмов і способи реагування наведено в табл.~\ref{tab:failures}.

\begin{table}[htbp]
\centering
\caption{Можливі збої та механізми відновлення}
\label{tab:failures}
\fontsize{11pt}{13pt}\selectfont
\setlength{\tabcolsep}{4pt}
\renewcommand{\arraystretch}{1.18}
\begin{tabularx}{\textwidth}{|>{\RaggedRight\arraybackslash}p{0.25\textwidth}|
                               >{\RaggedRight\arraybackslash}p{0.28\textwidth}|
                               Y|}
\hline
\textbf{Збій} & \textbf{Наслідок} & \textbf{Механізм відновлення} \\
\hline
Недостатньо коштів на \texttt{Node A} &
\texttt{Node A} повертає \texttt{ABORT}; зарахування не має бути зафіксоване. &
Координатор надсилає \texttt{ROLLBACK} усім підготовленим учасникам. \\
\hline
Неактивний рахунок на \texttt{Node B} &
\texttt{Node B} не може прийняти кошти. &
Для \texttt{Node A}, якщо він уже перебуває у стані \texttt{PREPARED}, виконується \texttt{ROLLBACK PREPARED}. \\
\hline
Розрив мережі до рішення \texttt{COMMIT} &
Координатор не отримує всі голоси. &
Координатор фіксує \texttt{ROLLBACK} і після відновлення з'єднання скасовує prepared transactions. \\
\hline
Збій після запису рішення \texttt{COMMIT} &
Один із вузлів може залишитися у стані \texttt{PREPARED}. &
Після відновлення координатор читає журнал рішення та повторює \texttt{COMMIT PREPARED}. \\
\hline
Збій координатора після \texttt{PREPARE} &
Блокування можуть утримуватися протягом тривалого часу. &
Відновлений координатор перевіряє \texttt{pg\_prepared\_xacts} і завершує операцію відповідно до журналу рішення. \\
\hline
\end{tabularx}
\end{table}

Отже, обрана модель дає змогу розмежувати локальне виконання операцій на вузлах і глобальне рішення координатора. Це забезпечує узгоджене завершення переказу навіть за наявності відмов, для яких у журналі координатора збережено остаточний стан операції.

\paragraphheading{Алгоритм ухвалення глобального рішення.}
Життєвий цикл переказу можна подати як послідовність станів глобальної операції. Спочатку координатор реєструє ідентифікатор транзакції й відкриває локальні операції на обох вузлах. Після виконання перевірок він отримує від учасників голоси. Якщо отримано два повідомлення \texttt{READY}, у журналі координатора фіксується рішення \texttt{COMMIT}; лише після цього команда передається вузлам. Якщо на будь-якому етапі до фіксації рішення отримано \texttt{ABORT} або вичерпано припустимий час очікування, координатор зберігає рішення \texttt{ROLLBACK}.

\begin{table}[htbp]
\centering
\caption{Стани глобальної операції переказу}
\label{tab:transaction-states}
\fontsize{11pt}{13pt}\selectfont
\setlength{\tabcolsep}{4pt}
\renewcommand{\arraystretch}{1.16}
\begin{tabularx}{\textwidth}{|>{\RaggedRight\arraybackslash}p{0.18\textwidth}|
                               >{\RaggedRight\arraybackslash}p{0.29\textwidth}|
                               Y|}
\hline
\textbf{Стан} & \textbf{Дія координатора} & \textbf{Очікуваний результат} \\
\hline
\texttt{STARTED} & Створення глобального ідентифікатора та відкриття локальних транзакцій. & Обидва вузли готові виконати перевірки й локальні зміни. \\
\hline
\texttt{PREPARING} & Надсилання запиту підготовки та збирання відповідей. & Кожний учасник повертає \texttt{READY} або \texttt{ABORT}. \\
\hline
\texttt{DECIDED} & Запис рішення \texttt{COMMIT} або \texttt{ROLLBACK} до журналу координатора. & Рішення може бути відновлене після збою координатора. \\
\hline
\texttt{FINISHED} & Надсилання завершальної команди всім учасникам. & Баланси зафіксовано узгоджено або повернуто до попередніх значень. \\
\hline
\end{tabularx}
\end{table}

Виконання завершальної фази має бути ідемпотентним з погляду координатора: після відновлення він повторно перевіряє наявність підготовлених транзакцій і надсилає те саме збережене рішення. Координатор не повинен самостійно замінювати раніше записаний \texttt{COMMIT} на \texttt{ROLLBACK}, оскільки це створило б відмінні результати на вузлах. Саме незмінність глобального рішення є умовою коректного відновлення розподіленої транзакції.
